\documentclass[12pt]{article}
\usepackage[UTF8]{inputenc}
\usepackage{xeCJK}
\setCJKmainfont{Sarabun}
\usepackage{enumitem}
\usepackage{geometry}
\geometry{a4paper, margin=2.5cm}
\usepackage{setspace}
\onehalfspacing

\begin{document}

\begin{center}
{\Large\textbf{ทบทวน (Review)}}\\[6pt]
{\normalsize วิชาฟิสิกส์ ม.ปลาย บทที่ 30}\\[4pt]
{\footnotesize ทั้งหมด 25 ข้อ}
\end{center}
\hrule
\bigskip

\section*{คำถาม in class}
\begin{enumerate}[label=\arabic*., itemsep=0.5\baselineskip, topsep=0.5\baselineskip, leftmargin=*]
\item วัตถุเคลื่อนที่ด้วยความเร่งคงที่ a จากหยุดนิ่ง ในเวลา t วินาที ขจัด s จะเป็นสมการใด?
\begin{enumerate}[label=\alph*)., itemsep=0.3\baselineskip, topsep=0.3\baselineskip]
\item ก\ 
\begin{minipage}[t]{0.9\linewidth}
\raggedright s = ½at²
\end{minipage}
\item ข\ 
\begin{minipage}[t]{0.9\linewidth}
\raggedright s = at²
\end{minipage}
\item ค\ 
\begin{minipage}[t]{0.9\linewidth}
\raggedright s = vt
\end{minipage}
\item ง\ 
\begin{minipage}[t]{0.9\linewidth}
\raggedright s = ½vt
\end{minipage}
\end{enumerate}
\vspace{-0.3\baselineskip}
\begin{small}
\textit{คำตอบ: ก}
\end{small}

\item พลังงานจลน์ของวัตถุมวล m เคลื่อนที่ด้วยความเร็ว v มีค่าเท่าใด?
\begin{enumerate}[label=\alph*)., itemsep=0.3\baselineskip, topsep=0.3\baselineskip]
\item ก\ 
\begin{minipage}[t]{0.9\linewidth}
\raggedright ½mv
\end{minipage}
\item ข\ 
\begin{minipage}[t]{0.9\linewidth}
\raggedright ½mv²
\end{minipage}
\item ค\ 
\begin{minipage}[t]{0.9\linewidth}
\raggedright mv²
\end{minipage}
\item ง\ 
\begin{minipage}[t]{0.9\linewidth}
\raggedright mgh
\end{minipage}
\end{enumerate}
\vspace{-0.3\baselineskip}
\begin{small}
\textit{คำตอบ: ข}
\end{small}

\item โมเมนตัม (Momentum) คำนวณได้จากสมการใด?
\begin{enumerate}[label=\alph*)., itemsep=0.3\baselineskip, topsep=0.3\baselineskip]
\item ก\ 
\begin{minipage}[t]{0.9\linewidth}
\raggedright p = mv
\end{minipage}
\item ข\ 
\begin{minipage}[t]{0.9\linewidth}
\raggedright p = FΔt
\end{minipage}
\item ค\ 
\begin{minipage}[t]{0.9\linewidth}
\raggedright p = ½mv²
\end{minipage}
\item ง\ 
\begin{minipage}[t]{0.9\linewidth}
\raggedright p = mgh
\end{minipage}
\end{enumerate}
\vspace{-0.3\baselineskip}
\begin{small}
\textit{คำตอบ: ก}
\end{small}

\item กฎของโอห์มระบุว่า V = IR โดยที่ V, I, R คืออะไร?
\begin{enumerate}[label=\alph*)., itemsep=0.3\baselineskip, topsep=0.3\baselineskip]
\item ก\ 
\begin{minipage}[t]{0.9\linewidth}
\raggedright V = กระแส, I = ความต่างศักย์, R = ความต้านทาน
\end{minipage}
\item ข\ 
\begin{minipage}[t]{0.9\linewidth}
\raggedright V = ความต่างศักย์, I = กระแส, R = ความต้านทาน
\end{minipage}
\item ค\ 
\begin{minipage}[t]{0.9\linewidth}
\raggedright V = กำลัง, I = ความต้านทาน, R = กระแส
\end{minipage}
\item ง\ 
\begin{minipage}[t]{0.9\linewidth}
\raggedright V = ความต่างศักย์, I = ความต้านทาน, R = กระแส
\end{minipage}
\end{enumerate}
\vspace{-0.3\baselineskip}
\begin{small}
\textit{คำตอบ: ข}
\end{small}

\item ความสัมพันธ์ระหว่างความยาวคลื่น λ, ความถี่ f และความเร็ว v คืออะไร?
\begin{enumerate}[label=\alph*)., itemsep=0.3\baselineskip, topsep=0.3\baselineskip]
\item ก\ 
\begin{minipage}[t]{0.9\linewidth}
\raggedright v = fλ
\end{minipage}
\item ข\ 
\begin{minipage}[t]{0.9\linewidth}
\raggedright v = f/λ
\end{minipage}
\item ค\ 
\begin{minipage}[t]{0.9\linewidth}
\raggedright v = λ/f
\end{minipage}
\item ง\ 
\begin{minipage}[t]{0.9\linewidth}
\raggedright v = f + λ
\end{minipage}
\end{enumerate}
\vspace{-0.3\baselineskip}
\begin{small}
\textit{คำตอบ: ก}
\end{small}

\item ปรากฏการณ์ใดพิสูจน์ว่าแสงมีสมบัติเป็นคลื่น?
\begin{enumerate}[label=\alph*)., itemsep=0.3\baselineskip, topsep=0.3\baselineskip]
\item ก\ 
\begin{minipage}[t]{0.9\linewidth}
\raggedright ปรากฏการณ์โฟโตอิเล็กตริก
\end{minipage}
\item ข\ 
\begin{minipage}[t]{0.9\linewidth}
\raggedright การเลี้ยวเบนของแสงผ่านเกรตติง
\end{minipage}
\item ค\ 
\begin{minipage}[t]{0.9\linewidth}
\raggedright การชนกระจก
\end{minipage}
\item ง\ 
\begin{minipage}[t]{0.9\linewidth}
\raggedright การสะท้อนกลับ
\end{minipage}
\end{enumerate}
\vspace{-0.3\baselineskip}
\begin{small}
\textit{คำตอบ: ข}
\end{small}

\item อิเล็กตรอนเปลี่ยนระดับพลังงานจาก n=3 ไป n=1 ของอะตอมไฮโดรเจน จะปล่อยรังสีที่มีความถี่เท่าใด? (Eₙ = -13.6/n² eV)
\begin{enumerate}[label=\alph*)., itemsep=0.3\baselineskip, topsep=0.3\baselineskip]
\item ก\ 
\begin{minipage}[t]{0.9\linewidth}
\raggedright ΔE = 12.09 eV
\end{minipage}
\item ข\ 
\begin{minipage}[t]{0.9\linewidth}
\raggedright ΔE = 13.6 eV
\end{minipage}
\item ค\ 
\begin{minipage}[t]{0.9\linewidth}
\raggedright ΔE = 1.51 eV
\end{minipage}
\item ง\ 
\begin{minipage}[t]{0.9\linewidth}
\raggedright ΔE = 10.2 eV
\end{minipage}
\end{enumerate}
\vspace{-0.3\baselineskip}
\begin{small}
\textit{คำตอบ: ก}
\end{small}

\item วัตถุถูกขว้างขึ้นไปในแนวดิ่งด้วยความเร็วต้น 20 m/s จะขึ้นไปได้สูงสุดกี่เมตร? (g = 10 m/s²)
\begin{enumerate}[label=\alph*)., itemsep=0.3\baselineskip, topsep=0.3\baselineskip]
\item ก\ 
\begin{minipage}[t]{0.9\linewidth}
\raggedright 10 m
\end{minipage}
\item ข\ 
\begin{minipage}[t]{0.9\linewidth}
\raggedright 20 m
\end{minipage}
\item ค\ 
\begin{minipage}[t]{0.9\linewidth}
\raggedright 40 m
\end{minipage}
\item ง\ 
\begin{minipage}[t]{0.9\linewidth}
\raggedright 4 m
\end{minipage}
\end{enumerate}
\vspace{-0.3\baselineskip}
\begin{small}
\textit{คำตอบ: ข}
\end{small}

\item ตัวต้านทาน 3 ตัวต่อขนานกัน มีความต้านทาน 6 Ω, 3 Ω และ 2 Ω ความต้านทานรวมเท่ากับ?
\begin{enumerate}[label=\alph*)., itemsep=0.3\baselineskip, topsep=0.3\baselineskip]
\item ก\ 
\begin{minipage}[t]{0.9\linewidth}
\raggedright 11 Ω
\end{minipage}
\item ข\ 
\begin{minipage}[t]{0.9\linewidth}
\raggedright 1 Ω
\end{minipage}
\item ค\ 
\begin{minipage}[t]{0.9\linewidth}
\raggedright 3 Ω
\end{minipage}
\item ง\ 
\begin{minipage}[t]{0.9\linewidth}
\raggedright 6 Ω
\end{minipage}
\end{enumerate}
\vspace{-0.3\baselineskip}
\begin{small}
\textit{คำตอบ: ข}
\end{small}

\item เงื่อนไขการแทรกสอดแบบเสริม (Constructive interference) คืออะไร?
\begin{enumerate}[label=\alph*)., itemsep=0.3\baselineskip, topsep=0.3\baselineskip]
\item ก\ 
\begin{minipage}[t]{0.9\linewidth}
\raggedright ΔL = (n + ½)λ
\end{minipage}
\item ข\ 
\begin{minipage}[t]{0.9\linewidth}
\raggedright ΔL = nλ
\end{minipage}
\item ค\ 
\begin{minipage}[t]{0.9\linewidth}
\raggedright ΔL = 2nλ
\end{minipage}
\item ง\ 
\begin{minipage}[t]{0.9\linewidth}
\raggedright ΔL = n/λ
\end{minipage}
\end{enumerate}
\vspace{-0.3\baselineskip}
\begin{small}
\textit{คำตอบ: ข}
\end{small}

\item สนามไฟฟ้าจากประจุจุด Q ที่ระยะ r มีขนาดเท่าใด? (k = 1/4πε₀)
\begin{enumerate}[label=\alph*)., itemsep=0.3\baselineskip, topsep=0.3\baselineskip]
\item ก\ 
\begin{minipage}[t]{0.9\linewidth}
\raggedright E = kQ/r²
\end{minipage}
\item ข\ 
\begin{minipage}[t]{0.9\linewidth}
\raggedright E = kQ/r
\end{minipage}
\item ค\ 
\begin{minipage}[t]{0.9\linewidth}
\raggedright E = kQ²/r²
\end{minipage}
\item ง\ 
\begin{minipage}[t]{0.9\linewidth}
\raggedright E = k/r²
\end{minipage}
\end{enumerate}
\vspace{-0.3\baselineskip}
\begin{small}
\textit{คำตอบ: ก}
\end{small}

\item ตามแบบจำลองอะตอมของโบร์ อิเล็กตรอนเปลี่ยนระดับพลังงานโดยการ...?
\begin{enumerate}[label=\alph*)., itemsep=0.3\baselineskip, topsep=0.3\baselineskip]
\item ก\ 
\begin{minipage}[t]{0.9\linewidth}
\raggedright ดูดกลืนโฟตอน
\end{minipage}
\item ข\ 
\begin{minipage}[t]{0.9\linewidth}
\raggedright ปล่อยหรือดูดกลืนโฟตอน
\end{minipage}
\item ค\ 
\begin{minipage}[t]{0.9\linewidth}
\raggedright ชนอะตอมอื่น
\end{minipage}
\item ง\ 
\begin{minipage}[t]{0.9\linewidth}
\raggedright เปลี่ยนโมเมนตัม
\end{minipage}
\end{enumerate}
\vspace{-0.3\baselineskip}
\begin{small}
\textit{คำตอบ: ข}
\end{small}

\item แรงสู่ศูนย์กลางสำหรับวัตถุมวล m เคลื่อนที่เป็นวงกลมรัศมี r ด้วยความเร็ว v มีค่าเท่าใด?
\begin{enumerate}[label=\alph*)., itemsep=0.3\baselineskip, topsep=0.3\baselineskip]
\item ก\ 
\begin{minipage}[t]{0.9\linewidth}
\raggedright F = mv/r
\end{minipage}
\item ข\ 
\begin{minipage}[t]{0.9\linewidth}
\raggedright F = mv²/r
\end{minipage}
\item ค\ 
\begin{minipage}[t]{0.9\linewidth}
\raggedright F = mvr
\end{minipage}
\item ง\ 
\begin{minipage}[t]{0.9\linewidth}
\raggedright F = mrv²
\end{minipage}
\end{enumerate}
\vspace{-0.3\baselineskip}
\begin{small}
\textit{คำตอบ: ข}
\end{small}

\item ความจุไฟฟ้า (Capacitance) ของตัวเก็บประจุแบบแผ่นขนาน C = ε₀A/d โดย A, d, ε₀ คืออะไร?
\begin{enumerate}[label=\alph*)., itemsep=0.3\baselineskip, topsep=0.3\baselineskip]
\item ก\ 
\begin{minipage}[t]{0.9\linewidth}
\raggedright A = พื้นที่, d = ระยะห่าง, ε₀ = ค่าคงที่ไฟฟ้า
\end{minipage}
\item ข\ 
\begin{minipage}[t]{0.9\linewidth}
\raggedright A = ประจุ, d = ความต่างศักย์, ε₀ = สนามไฟฟ้า
\end{minipage}
\item ค\ 
\begin{minipage}[t]{0.9\linewidth}
\raggedright A = ระยะห่าง, d = พื้นที่, ε₀ = ค่าคงที่ไฟฟ้า
\end{minipage}
\item ง\ 
\begin{minipage}[t]{0.9\linewidth}
\raggedright A = ความต่างศักย์, d = ความจุ, ε₀ = พื้นที่
\end{minipage}
\end{enumerate}
\vspace{-0.3\baselineskip}
\begin{small}
\textit{คำตอบ: ก}
\end{small}

\item แหล่งกำเนิดเสียงเคลื่อนที่เข้าหาผู้สังเกตที่หยุดนิ่ง ผู้สังเกตจะได้ยินเสียงความถี่อย่างไร?
\begin{enumerate}[label=\alph*)., itemsep=0.3\baselineskip, topsep=0.3\baselineskip]
\item ก\ 
\begin{minipage}[t]{0.9\linewidth}
\raggedright ต่ำลง
\end{minipage}
\item ข\ 
\begin{minipage}[t]{0.9\linewidth}
\raggedright สูงขึ้น
\end{minipage}
\item ค\ 
\begin{minipage}[t]{0.9\linewidth}
\raggedright เท่าเดิม
\end{minipage}
\item ง\ 
\begin{minipage}[t]{0.9\linewidth}
\raggedright ขึ้นอยู่กับทิศทางลม
\end{minipage}
\end{enumerate}
\vspace{-0.3\baselineskip}
\begin{small}
\textit{คำตอบ: ข}
\end{small}

\item นิวเคลียสที่เสถียรที่สุดในตารางนิวไคลด์คือธาตุใด?
\begin{enumerate}[label=\alph*)., itemsep=0.3\baselineskip, topsep=0.3\baselineskip]
\item ก\ 
\begin{minipage}[t]{0.9\linewidth}
\raggedright ไฮโดรเจน-1
\end{minipage}
\item ข\ 
\begin{minipage}[t]{0.9\linewidth}
\raggedright ยูเรเนียม-235
\end{minipage}
\item ค\ 
\begin{minipage}[t]{0.9\linewidth}
\raggedright เหล็ก-56
\end{minipage}
\item ง\ 
\begin{minipage}[t]{0.9\linewidth}
\raggedright ทองคำ-197
\end{minipage}
\end{enumerate}
\vspace{-0.3\baselineskip}
\begin{small}
\textit{คำตอบ: ค}
\end{small}

\item วัตถุเคลื่อนที่จากจุด A ไป B ด้วยความเร็ว 4 m/s ในเวลา 5 s แล้วกลับจาก B ไป A ด้วยความเร็ว 2 m/s ความเร็วเฉลี่ยทั้งเที่ยวคือเท่าใด?
\begin{enumerate}[label=\alph*)., itemsep=0.3\baselineskip, topsep=0.3\baselineskip]
\item ก\ 
\begin{minipage}[t]{0.9\linewidth}
\raggedright 3 m/s
\end{minipage}
\item ข\ 
\begin{minipage}[t]{0.9\linewidth}
\raggedright 2.67 m/s
\end{minipage}
\item ค\ 
\begin{minipage}[t]{0.9\linewidth}
\raggedright 3.5 m/s
\end{minipage}
\item ง\ 
\begin{minipage}[t]{0.9\linewidth}
\raggedright 4 m/s
\end{minipage}
\end{enumerate}
\vspace{-0.3\baselineskip}
\begin{small}
\textit{คำตอบ: ข}
\end{small}

\item แรงที่กระทำต่อประจุไฟฟ้า q เคลื่อนที่ด้วยความเร็ว v ในสนามแม่เหล็ก B มีค่าเท่าใด?
\begin{enumerate}[label=\alph*)., itemsep=0.3\baselineskip, topsep=0.3\baselineskip]
\item ก\ 
\begin{minipage}[t]{0.9\linewidth}
\raggedright F = qvB
\end{minipage}
\item ข\ 
\begin{minipage}[t]{0.9\linewidth}
\raggedright F = qvB sinθ
\end{minipage}
\item ค\ 
\begin{minipage}[t]{0.9\linewidth}
\raggedright F = qv²B
\end{minipage}
\item ง\ 
\begin{minipage}[t]{0.9\linewidth}
\raggedright F = qB/v
\end{minipage}
\end{enumerate}
\vspace{-0.3\baselineskip}
\begin{small}
\textit{คำตอบ: ข}
\end{small}

\item ความเข้มเสียง (Sound intensity) ที่ระยะ r จากแหล่งกำเนิดกำลัง P มีค่าเท่าใด?
\begin{enumerate}[label=\alph*)., itemsep=0.3\baselineskip, topsep=0.3\baselineskip]
\item ก\ 
\begin{minipage}[t]{0.9\linewidth}
\raggedright I = P/4πr²
\end{minipage}
\item ข\ 
\begin{minipage}[t]{0.9\linewidth}
\raggedright I = P/r²
\end{minipage}
\item ค\ 
\begin{minipage}[t]{0.9\linewidth}
\raggedright I = 4πr²/P
\end{minipage}
\item ง\ 
\begin{minipage}[t]{0.9\linewidth}
\raggedright I = P/4πr
\end{minipage}
\end{enumerate}
\vspace{-0.3\baselineskip}
\begin{small}
\textit{คำตอบ: ก}
\end{small}

\item สารกัมมันตรังสีมีค่าครึ่งชีวิต 10 ปี ถ้าเริ่มต้นมีมวล 80 g จะเหลือเท่าใดหลังผ่านไป 30 ปี?
\begin{enumerate}[label=\alph*)., itemsep=0.3\baselineskip, topsep=0.3\baselineskip]
\item ก\ 
\begin{minipage}[t]{0.9\linewidth}
\raggedright 10 g
\end{minipage}
\item ข\ 
\begin{minipage}[t]{0.9\linewidth}
\raggedright 5 g
\end{minipage}
\item ค\ 
\begin{minipage}[t]{0.9\linewidth}
\raggedright 2.5 g
\end{minipage}
\item ง\ 
\begin{minipage}[t]{0.9\linewidth}
\raggedright 1 g
\end{minipage}
\end{enumerate}
\vspace{-0.3\baselineskip}
\begin{small}
\textit{คำตอบ: ก}
\end{small}

\item งาน (Work) ที่ทำโดยแรง F ที่ทำมุม θ กับการเคลื่อนที่ระยะ s คืออะไร?
\begin{enumerate}[label=\alph*)., itemsep=0.3\baselineskip, topsep=0.3\baselineskip]
\item ก\ 
\begin{minipage}[t]{0.9\linewidth}
\raggedright W = Fs sinθ
\end{minipage}
\item ข\ 
\begin{minipage}[t]{0.9\linewidth}
\raggedright W = Fs cosθ
\end{minipage}
\item ค\ 
\begin{minipage}[t]{0.9\linewidth}
\raggedright W = Fs tanθ
\end{minipage}
\item ง\ 
\begin{minipage}[t]{0.9\linewidth}
\raggedright W = Fs
\end{minipage}
\end{enumerate}
\vspace{-0.3\baselineskip}
\begin{small}
\textit{คำตอบ: ข}
\end{small}

\item กำลังไฟฟ้า P = IV = I²R = V²/R สูตรใดถูกต้องสำหรับตัวต้านทาน R?
\begin{enumerate}[label=\alph*)., itemsep=0.3\baselineskip, topsep=0.3\baselineskip]
\item ก\ 
\begin{minipage}[t]{0.9\linewidth}
\raggedright ทุกสูตรถูกต้อง
\end{minipage}
\item ข\ 
\begin{minipage}[t]{0.9\linewidth}
\raggedright P = I²R เท่านั้น
\end{minipage}
\item ค\ 
\begin{minipage}[t]{0.9\linewidth}
\raggedright P = V²/R เท่านั้น
\end{minipage}
\item ง\ 
\begin{minipage}[t]{0.9\linewidth}
\raggedright P = IV เท่านั้น
\end{minipage}
\end{enumerate}
\vspace{-0.3\baselineskip}
\begin{small}
\textit{คำตอบ: ก}
\end{small}

\item สมการคลื่นแม่เหล็กไฟฟ้าของแสง E = hf หมายความว่าอะไร?
\begin{enumerate}[label=\alph*)., itemsep=0.3\baselineskip, topsep=0.3\baselineskip]
\item ก\ 
\begin{minipage}[t]{0.9\linewidth}
\raggedright แสงเป็นได้ทั้งคลื่นและอนุภาค
\end{minipage}
\item ข\ 
\begin{minipage}[t]{0.9\linewidth}
\raggedright แสงมีพลังงานเป็นสัดส่วนโดยตรงกับความถี่
\end{minipage}
\item ค\ 
\begin{minipage}[t]{0.9\linewidth}
\raggedright ความถี่เป็นสัดส่วนกับความยาวคลื่น
\end{minipage}
\item ง\ 
\begin{minipage}[t]{0.9\linewidth}
\raggedright ถูกทั้ง ก และ ข
\end{minipage}
\end{enumerate}
\vspace{-0.3\baselineskip}
\begin{small}
\textit{คำตอบ: ง}
\end{small}

\item ปฏิกิริยานิวเคลียร์ประเภทใดที่เกิดในดวงอาทิตย์?
\begin{enumerate}[label=\alph*)., itemsep=0.3\baselineskip, topsep=0.3\baselineskip]
\item ก\ 
\begin{minipage}[t]{0.9\linewidth}
\raggedright ฟิชชัน
\end{minipage}
\item ข\ 
\begin{minipage}[t]{0.9\linewidth}
\raggedright ฟิวชัน (Proton-Proton chain)
\end{minipage}
\item ค\ 
\begin{minipage}[t]{0.9\linewidth}
\raggedright การสลายตัวของคาร์บอน
\end{minipage}
\item ง\ 
\begin{minipage}[t]{0.9\linewidth}
\raggedright การเปลี่ยนสถานะของสสาร
\end{minipage}
\end{enumerate}
\vspace{-0.3\baselineskip}
\begin{small}
\textit{คำตอบ: ข}
\end{small}

\item ข้อสอบ PAT3 มักให้ค่าคงที่ต่างๆ ในตาราง ถ้าโจทย์ให้มวลอะตอมของอนุภาค แล้วถามหาพลังงานยึดเหนี่ยว ขั้นตอนแรกที่ควรทำคืออะไร?
\begin{enumerate}[label=\alph*)., itemsep=0.3\baselineskip, topsep=0.3\baselineskip]
\item ก\ 
\begin{minipage}[t]{0.9\linewidth}
\raggedright คูณมวลกับ g
\end{minipage}
\item ข\ 
\begin{minipage}[t]{0.9\linewidth}
\raggedright หามวลต่ำ (Mass defect) ก่อน แล้วใช้ E = Δmc²
\end{minipage}
\item ค\ 
\begin{minipage}[t]{0.9\linewidth}
\raggedright ใช้สูตร E = mc² โดยตรง
\end{minipage}
\item ง\ 
\begin{minipage}[t]{0.9\linewidth}
\raggedright หาพลังงานจลน์
\end{minipage}
\end{enumerate}
\vspace{-0.3\baselineskip}
\begin{small}
\textit{คำตอบ: ข}
\end{small}

\end{enumerate}
\end{document}